\documentclass[12pt,a4paper]{ctexart}
\usepackage[margin=2.2cm]{geometry}
\usepackage{graphicx}
\usepackage{booktabs}
\usepackage{amsmath}
\usepackage{float}
\usepackage{hyperref}
\usepackage{xcolor}
\usepackage{listings}

\lstset{
  basicstyle=\ttfamily\small,
  keywordstyle=\color{blue!70!black},
  commentstyle=\color{green!40!black},
  stringstyle=\color{orange!60!black},
  showstringspaces=false,
  breaklines=true,
  frame=single,
  tabsize=4
}

\newcommand{\reporttitle}{Cache 基本性能分析}
\newcommand{\authname}{姓名}
\newcommand{\authid}{学号}

\title{\reporttitle}
\author{\authname：\underline{李祥熙} \quad \authid：\underline{2234412799}}
\date{\today}

\begin{document}
\maketitle
\tableofcontents
\newpage

\section{实验目的}
本实验的目标如下：
\begin{itemize}
  \item 加深对 Cache 大小、相联度、块大小、替换算法等因素对 Cache 性能影响的理解；
  \item 自行设计并实现一个可配置的 Cache 模拟器（模拟器 A）；
  \item 使用真实程序的访存地址流（trace）对模拟器进行测试；
  \item 通过控制变量的方式，分别考察各个参数对失效率的影响，并给出分析结论。
\end{itemize}

\section{实验过程}
\subsection{实验环境}
\begin{itemize}
  \item 操作系统：macOS（Darwin 25.5.0）
  \item 编程语言：C++（C++17 标准）
  \item 编译器：g++ \texttt{-O2}
  \item 辅助工具：Bash 脚本（批量实验），awk（数据整理）
\end{itemize}
本实验选择模拟器 A，即完全自行实现 Cache 模拟器。选择 C++ 是因为每个 trace 文件约有
100 万条记录，而完整实验需要遍历 256 种参数组合，C++ 的执行效率可以保证整套实验在数十秒内完成。

\subsection{数据集与预处理}
实验使用本次实验提供的 4 个 trace 文件，它们是 SPEC 基准程序的访存地址流：
\begin{itemize}
  \item \texttt{022.li.din}：LISP 解释器，约 100 万条记录；
  \item \texttt{047.tomcatv.din}：向量化的网格生成程序，访存以数组为主；
  \item \texttt{078.swm256.din}：浅水方程数值模拟，以读为主；
  \item \texttt{085.gcc.din}：GCC 编译器，访存模式不规则。
\end{itemize}
每个 trace 文件中，每一行的格式为 \texttt{access\_type address [size/data]}，其中
\texttt{access\_type} 取值含义为：\texttt{0} 表示 load data（读），\texttt{1} 表示
store data（写），\texttt{2} 表示 fetch instruction（取指）；\texttt{address} 为 16
进制表示的 32 位地址。一个典型的片段如下：
\begin{lstlisting}
2 400190 8fa40000
0 7ffebc64 0
2 400194 3c1c1001
1 100039b8 0
\end{lstlisting}
按照实验要求 H（“为简化模拟器实现，仅考虑 load data 和 store data”），模拟器在读取地址流时
\textbf{只处理类型 0 和类型 1 的访存记录，忽略类型 2 的取指记录}。每次访存按 4 字节处理。

\subsection{实现步骤}
\begin{enumerate}
  \item 设计 Cache 的数据结构：将 Cache 组织为“组—路—行”的三级结构，每行只保存 tag 和有效位；
  \item 实现地址划分逻辑：根据块大小和组数，把访存地址拆分为 tag、组索引、块内偏移三部分；
  \item 实现 LRU 替换算法：每组内的若干路按最近使用顺序排列，命中或换入时把该行移到最前；
  \item 实现写分配（write-allocate）策略：写失效时也把数据块调入 Cache；
  \item 实现命令行参数解析，支持 \texttt{trace\_input\_file}、\texttt{cache\_size}、
        \texttt{cache\_associativity}、\texttt{cache\_block\_size} 四个参数；
  \item 统计并输出读/写访存数、读/写失效数、替换块数以及各项失效率；
  \item 编写批处理脚本，遍历全部参数组合，把结果汇总成 CSV 文件用于分析。
\end{enumerate}

\section{实验内容}
\subsection{模拟器设计思想}
模拟器采用经典的\textbf{组相联 Cache} 模型，这是直接映射（相联度为 1）和全相联（组数为 1）
的统一形式。其核心参数关系为：
\begin{align}
\text{块数} &= \frac{\text{cache\_size}}{\text{block\_size}} \\
\text{组数} &= \frac{\text{块数}}{\text{cache\_associativity}}
\end{align}

对于一个访存地址，去掉低位的\textbf{块内偏移}后得到块地址；块地址的低位作为\textbf{组索引}，
用于定位到具体某一组；剩余的高位作为 \textbf{tag}，用于在组内比较是否命中。由于
\texttt{cache\_size}、\texttt{block\_size}、\texttt{assoc} 均为 2 的幂，组数也必然是 2 的幂，
因此组索引可以用一次按位与（\texttt{\& (numSets-1)}）高效计算。地址划分的核心代码如下：
\begin{lstlisting}
uint64_t blockAddr = addr >> offsetBits_;
uint64_t index = numSets_ > 1 ? (blockAddr & (numSets_ - 1)) : 0;
uint64_t tag   = blockAddr >> indexBits_;
\end{lstlisting}

\textbf{LRU 替换算法}的实现采用了一个简洁的技巧：把每一组的各路保存在一个向量中，
\textbf{始终让最近使用的行排在最前面（front = MRU），最久未使用的行排在最后（back = LRU）}。
这样，命中时只需把命中行移到队首；失效换入时，若该组已满（最后一行有效）则发生一次替换，
直接弹出队尾、把新行插入队首即可。由于相联度最多为 8，这种线性操作的开销可以忽略。
单次访存处理的核心逻辑如下：
\begin{lstlisting}
auto& set = sets_[index];
for (size_t i = 0; i < set.size(); ++i) {
    if (set[i].valid && set[i].tag == tag) {   // ---- 命中 ----
        moveToFront(set, i);                    // 更新 LRU 顺序
        return;
    }
}
// ---- 失效 ----（写分配：失效总是把块调入 Cache）
if (isWrite) ++writeMisses_; else ++readMisses_;
if (set.back().valid) ++replacements_;          // 该组已满，发生替换
set.pop_back();
set.insert(set.begin(), Line{tag, true});       // 新块插入队首
\end{lstlisting}

\subsection{模拟器特色}
\begin{itemize}
  \item \textbf{统一的组相联模型}：相联度为 1 时即直接映射，组数为 1 时即全相联，一套代码覆盖全部情况；
  \item \textbf{参数合法性检查}：要求三个参数均为 2 的幂，并校验 \texttt{block\_size} $\le$
        \texttt{cache\_size} 以及 \texttt{(cache\_size/block\_size) \% assoc == 0}，避免非法配置；
  \item \textbf{友好的参数输入}：大小参数支持 \texttt{K}/\texttt{M} 后缀，例如 \texttt{--cache\_size 32K}；
  \item \textbf{两种输出模式}：默认输出可读的统计表格，加 \texttt{--csv} 则输出单行机器可读数据，便于批量汇总；
  \item \textbf{流式读取}：逐行处理 trace，内存占用与文件大小无关。
\end{itemize}

\subsection{实验参数组合}
按照实验要求，固定使用 \textbf{LRU 替换算法}和\textbf{写分配}策略，对以下参数进行充分测试：
\begin{table}[H]
\centering
\caption{实验参数取值范围}
\label{tab:params}
\begin{tabular}{ll}
\toprule
参数 & 取值 \\
\midrule
trace 文件 & 022.li、047.tomcatv、078.swm256、085.gcc \\
Cache 大小 & 8KB、16KB、32KB、64KB \\
相联度 & 1、2、4、8 \\
块大小 & 16B、32B、64B、128B \\
\bottomrule
\end{tabular}
\end{table}
四个维度各取 4 个值，共 $4\times4\times4\times4 = 256$ 种组合，全部测试结果汇总在
\texttt{results/results.csv} 中。下文通过\textbf{控制变量法}，每次只改变一个参数来观察其影响。

\section{实验结果}
\subsection{各 trace 的访存特征}
首先给出各 trace 在仅统计 load/store 后的访存规模，这有助于理解后面失效率的差异。
\begin{table}[H]
\centering
\caption{各 trace 的访存特征（仅 load/store）}
\label{tab:profile}
\begin{tabular}{lrrr}
\toprule
trace & 读次数 & 写次数 & 写占比 \\
\midrule
022.li      & 155399 & 102350 & 39.7\% \\
047.tomcatv & 253723 & 130733 & 34.0\% \\
078.swm256  & 192459 &   9894 &  4.9\% \\
085.gcc     & 139816 &  80669 & 36.6\% \\
\bottomrule
\end{tabular}
\end{table}
可以看到 078.swm256 的写操作占比极低（仅 4.9\%），是典型的“读密集”程序；其余三个程序读写
比例较为均衡。这一特征会直接影响它们对写失效的敏感程度。

\subsection{某一具体配置的详细统计}
为说明模拟器输出的完整内容，下面给出在 \textbf{Cache=32KB、相联度=4、块大小=64B} 这一组配置下，
对 022.li 运行得到的原始输出（即默认表格模式）：
\begin{lstlisting}
==================== Cache Simulation Result ====================
 Trace file       : trace files/022.li.din
 Cache size       : 32768 B (32 KB)
 Block size       : 64 B
 Associativity    : 4-way
 Number of sets   : 128
 Replacement      : LRU      Write-miss: write-allocate
-----------------------------------------------------------------
 Data accesses    : 257749  (read 155399, write 102350)
 Read  misses     : 452
 Write misses     : 872
 Replaced blocks  : 812
 Read  miss rate  : 0.2909 %
 Write miss rate  : 0.8520 %
 Total miss rate  : 0.5137 %
=================================================================
\end{lstlisting}
该输出包含了实验要求的全部统计项：读/写访存数、读/写失效数、替换块数，以及读、写和总的失效率。
其中“Number of sets = 128”由 $32768 / 64 / 4 = 128$ 计算得到；“Replaced blocks”小于失效总数
（$452+872=1324$），差值 $1324-812=512$ 正是各组首次被填满之前的\textbf{冷启动失效}（强制失效），
它们不引起替换。

把该配置推广到全部 4 个 trace，得到表~\ref{tab:detail} 的横向对比：
\begin{table}[H]
\centering
\caption{Cache=32KB、相联度=4、块大小=64B 下的详细统计}
\label{tab:detail}
\begin{tabular}{lrrrrr}
\toprule
trace & 读失效 & 写失效 & 替换块数 & 读失效率 & 写失效率 \\
\midrule
022.li      & 452 &  872 &  812 & 0.29\% & 0.85\% \\
047.tomcatv & 168 & 8573 & 8229 & 0.07\% & 6.56\% \\
078.swm256  &  33 &  526 &   57 & 0.02\% & 5.32\% \\
085.gcc     & 683 & 1330 & 1501 & 0.49\% & 1.65\% \\
\bottomrule
\end{tabular}
\end{table}
表~\ref{tab:detail} 揭示了一个有趣现象：047.tomcatv 和 078.swm256 的\textbf{读失效率极低
但写失效率较高}。这是因为它们是数值计算程序，对数组的读取具有很好的空间局部性（块被读入后被反复读取），
而写操作往往是“先写后才会再用”的初始化式写入，在写分配策略下每次写新块都会产生一次写失效。

\subsection{Cache 大小对失效率的影响}
固定\textbf{相联度=4、块大小=64B}，改变 Cache 大小，总失效率（\%）如表~\ref{tab:size} 所示。
\begin{table}[H]
\centering
\caption{Cache 大小对总失效率的影响（相联度=4，块大小=64B，单位：\%）}
\label{tab:size}
\begin{tabular}{lrrrr}
\toprule
trace & 8KB & 16KB & 32KB & 64KB \\
\midrule
022.li      & 1.4180 & 0.7845 & 0.5137 & 0.3845 \\
047.tomcatv & 2.6461 & 2.2744 & 2.2736 & 2.2653 \\
078.swm256  & 0.2762 & 0.2762 & 0.2762 & 0.2762 \\
085.gcc     & 2.8696 & 1.2055 & 0.9130 & 0.7597 \\
\bottomrule
\end{tabular}
\end{table}
结果分析：
\begin{itemize}
  \item \textbf{整体趋势是 Cache 越大失效率越低}，这符合理论预期：容量增大减少了容量失效；
  \item 022.li 和 085.gcc 受容量影响最明显（085.gcc 从 8KB 的 2.87\% 降到 64KB 的 0.76\%），
        说明它们的工作集较大、对容量较敏感；
  \item 047.tomcatv 在 16KB 之后几乎不再下降，078.swm256 则\textbf{从头到尾保持 0.2762\% 不变}。
        这是因为它们剩下的失效几乎都是\textbf{强制失效}（每个数据块第一次访问必然失效），
        而强制失效与 Cache 大小无关，因此增大容量无济于事。
\end{itemize}

\subsection{相联度对失效率的影响}
固定 \textbf{Cache=16KB、块大小=64B}，改变相联度，总失效率（\%）如表~\ref{tab:assoc} 所示。
\begin{table}[H]
\centering
\caption{相联度对总失效率的影响（Cache=16KB，块大小=64B，单位：\%）}
\label{tab:assoc}
\begin{tabular}{lrrrr}
\toprule
trace & 1 路 & 2 路 & 4 路 & 8 路 \\
\midrule
022.li      & 2.2184 & 0.9432 & 0.7845 & 0.7333 \\
047.tomcatv & 3.0050 & 2.3815 & 2.2744 & 2.2827 \\
078.swm256  & 0.3889 & 0.2762 & 0.2762 & 0.2762 \\
085.gcc     & 3.1671 & 1.5248 & 1.2055 & 1.0985 \\
\bottomrule
\end{tabular}
\end{table}
结果分析：
\begin{itemize}
  \item \textbf{提高相联度能降低失效率}，因为它减少了多个地址映射到同一组而互相驱逐造成的\textbf{冲突失效}；
  \item \textbf{收益主要集中在从 1 路到 2 路的阶段}：022.li 从 2.22\% 骤降到 0.94\%，
        085.gcc 从 3.17\% 降到 1.52\%，几乎减半；
  \item 而从 4 路提高到 8 路时改善很小，047.tomcatv 甚至略有回升（2.2744\%→2.2827\%）。
        这说明当冲突失效已被基本消除后，继续提高相联度的边际收益递减，剩下的失效以强制失效为主；
  \item 这一规律与课堂上的结论一致：\textbf{“2 路组相联 Cache 大致相当于容量翻倍的直接映射 Cache”}，
        即提高相联度和增大容量在降低冲突失效上有相似效果。
\end{itemize}

\subsection{块大小对失效率的影响}
固定 \textbf{Cache=32KB、相联度=4}，改变块大小，总失效率（\%）如表~\ref{tab:block} 所示。
\begin{table}[H]
\centering
\caption{块大小对总失效率的影响（Cache=32KB，相联度=4，单位：\%）}
\label{tab:block}
\begin{tabular}{lrrrr}
\toprule
trace & 16B & 32B & 64B & 128B \\
\midrule
022.li      & 1.6729 & 0.9327 & 0.5137 & 0.2863 \\
047.tomcatv & 8.6767 & 4.4073 & 2.2736 & 1.2051 \\
078.swm256  & 1.0546 & 0.5387 & 0.2762 & 0.1428 \\
085.gcc     & 2.6437 & 1.5058 & 0.9130 & 0.5710 \\
\bottomrule
\end{tabular}
\end{table}
结果分析：
\begin{itemize}
  \item \textbf{在本实验的取值范围内，增大块大小对降低失效率效果非常显著}，几乎每翻一倍块大小，失效率就大致减半；
  \item 这是因为更大的块一次调入更多相邻数据，充分利用了程序的\textbf{空间局部性}，从而减少了强制失效的次数；
  \item 数值计算程序 047.tomcatv 的改善最为剧烈（从 16B 的 8.68\% 降到 128B 的 1.21\%），
        因为它顺序访问数组，空间局部性极强，大块能一次性覆盖更多后续访问；
  \item 需要说明的是，块大小并非越大越好：块过大时会减少 Cache 中的总块数，可能引入更多冲突失效，
        并增大每次失效的传输开销。本实验最大块为 128B，尚未触及这一拐点，所以表现为单调下降。
\end{itemize}

\section{总结}
本实验自行设计并实现了一个组相联 Cache 模拟器（模拟器 A），采用 LRU 替换和写分配策略，
支持 \texttt{trace\_input\_file}、\texttt{cache\_size}、\texttt{cache\_associativity}、
\texttt{cache\_block\_size} 四个参数，并用 4 个真实程序的 trace 进行了 256 组组合测试。
主要完成内容与结论如下：
\begin{itemize}
  \item 实现了正确、高效的 Cache 模拟器，一套代码统一覆盖直接映射、组相联和全相联；
  \item \textbf{增大 Cache 容量}主要减少容量失效，对工作集大的程序（gcc、li）效果好，
        但无法消除强制失效；
  \item \textbf{提高相联度}主要减少冲突失效，收益集中在 1 路到 2 路阶段，之后边际递减；
  \item \textbf{增大块大小}通过利用空间局部性显著减少强制失效，对顺序访存的数值程序（tomcatv）效果最突出；
  \item 三类参数的效果与三种失效（强制、容量、冲突）的成因一一对应，实验数据与课堂理论高度吻合。
\end{itemize}

\textbf{实验感悟}：通过亲手实现模拟器，我对 Cache 的工作原理有了远比看公式更直观的理解。
尤其是观察到“增大容量对某些程序毫无作用、而增大块大小却效果惊人”这类现象时，
真正体会到 Cache 性能优化必须结合具体程序的访存模式（局部性）来权衡，
不存在对所有程序都最优的单一配置。控制变量法在分析多参数系统时也非常有效，
它让每个参数的独立影响清晰可见。

\section*{附录：实验源代码}
\subsection*{A.1 主要源码：src/cache\_sim.cpp}
\begin{lstlisting}
// 配置化的组相联 Cache 模拟器（模拟器 A）
//   替换策略：LRU      写失效策略：写分配（write-allocate）
//   仅模拟 load(0) 和 store(1)，忽略取指(2)
#include <cstdint>
#include <cstdio>
#include <cstdlib>
#include <string>
#include <vector>
#include <fstream>

struct Line { uint64_t tag = 0; bool valid = false; };

class Cache {
public:
    Cache(uint64_t cacheSize, uint32_t blockSize, uint32_t assoc)
        : blockSize_(blockSize), assoc_(assoc) {
        numBlocks_ = cacheSize / blockSize;
        numSets_   = numBlocks_ / assoc;
        offsetBits_ = log2u(blockSize);
        indexBits_  = log2u(numSets_);
        sets_.assign(numSets_, std::vector<Line>(assoc));  // front=MRU
    }
    void access(uint64_t addr, bool isWrite) {
        uint64_t blockAddr = addr >> offsetBits_;
        uint64_t index = numSets_ > 1 ? (blockAddr & (numSets_ - 1)) : 0;
        uint64_t tag   = blockAddr >> indexBits_;
        if (isWrite) ++writes_; else ++reads_;
        auto& set = sets_[index];
        for (size_t i = 0; i < set.size(); ++i) {
            if (set[i].valid && set[i].tag == tag) {   // 命中
                moveToFront(set, i); return;
            }
        }
        if (isWrite) ++writeMisses_; else ++readMisses_; // 失效
        if (set.back().valid) ++replacements_;           // 满组->替换
        set.pop_back();
        set.insert(set.begin(), Line{tag, true});        // 写分配
    }
    uint64_t reads() const { return reads_; }
    uint64_t writes() const { return writes_; }
    uint64_t readMisses() const { return readMisses_; }
    uint64_t writeMisses() const { return writeMisses_; }
    uint64_t replacements() const { return replacements_; }
    uint64_t numSets() const { return numSets_; }
private:
    static uint32_t log2u(uint64_t x) { uint32_t b=0; while(x>1){x>>=1;++b;} return b; }
    static void moveToFront(std::vector<Line>& set, size_t i) {
        Line tmp = set[i]; set.erase(set.begin()+i); set.insert(set.begin(), tmp);
    }
    uint32_t blockSize_, assoc_;
    uint64_t numBlocks_, numSets_;
    uint32_t offsetBits_, indexBits_;
    std::vector<std::vector<Line>> sets_;
    uint64_t reads_=0, writes_=0, readMisses_=0, writeMisses_=0, replacements_=0;
};

static bool isPow2(uint64_t x){ return x && !(x&(x-1)); }
static uint64_t parseSize(const std::string& s){     // 支持 K/M 后缀
    char suf = s.empty()?0:s.back(); uint64_t mul=1; std::string num=s;
    if(suf=='K'||suf=='k'){mul=1024ULL;num.pop_back();}
    else if(suf=='M'||suf=='m'){mul=1024ULL*1024;num.pop_back();}
    return std::strtoull(num.c_str(),nullptr,10)*mul;
}

int main(int argc, char** argv){
    std::string tracePath; uint64_t cacheSize=0, blockSize=0; uint32_t assoc=0;
    bool csv=false;
    for(int i=1;i<argc;++i){
        std::string a=argv[i];
        auto next=[&]()->std::string{ if(i+1>=argc) std::exit(1); return argv[++i]; };
        if(a=="--trace") tracePath=next();
        else if(a=="--cache_size") cacheSize=parseSize(next());
        else if(a=="--assoc") assoc=(uint32_t)std::strtoul(next().c_str(),nullptr,10);
        else if(a=="--block_size") blockSize=parseSize(next());
        else if(a=="--csv") csv=true;
    }
    if(tracePath.empty()||!cacheSize||!blockSize||!assoc) return 1;
    if(!isPow2(cacheSize)||!isPow2(blockSize)||!isPow2(assoc)) return 1;
    if(blockSize>cacheSize||(cacheSize/blockSize)%assoc!=0) return 1;

    std::ifstream in(tracePath); if(!in) return 1;
    Cache cache(cacheSize,(uint32_t)blockSize,assoc);
    int type; std::string addrHex, rest;
    while(in >> type >> addrHex){
        std::getline(in, rest);                  // 丢弃 size/data 字段
        if(type==0)      cache.access(std::strtoull(addrHex.c_str(),nullptr,16), false);
        else if(type==1) cache.access(std::strtoull(addrHex.c_str(),nullptr,16), true);
        // type==2 取指：忽略
    }
    uint64_t reads=cache.reads(), writes=cache.writes();
    uint64_t rMiss=cache.readMisses(), wMiss=cache.writeMisses();
    uint64_t accesses=reads+writes, misses=rMiss+wMiss;
    auto rate=[](uint64_t m,uint64_t n){ return n?100.0*(double)m/(double)n:0.0; };
    if(csv){
        std::printf("%s,%llu,%u,%llu,%llu,%llu,%llu,%llu,%llu,%.4f,%.4f,%.4f\n",
            tracePath.c_str(),(unsigned long long)cacheSize,assoc,(unsigned long long)blockSize,
            (unsigned long long)reads,(unsigned long long)writes,
            (unsigned long long)rMiss,(unsigned long long)wMiss,
            (unsigned long long)cache.replacements(),
            rate(rMiss,reads),rate(wMiss,writes),rate(misses,accesses));
    } else {
        std::printf(" Data accesses : %llu (read %llu, write %llu)\n",
            (unsigned long long)accesses,(unsigned long long)reads,(unsigned long long)writes);
        std::printf(" Read miss %.4f%%  Write miss %.4f%%  Total %.4f%%\n",
            rate(rMiss,reads),rate(wMiss,writes),rate(misses,accesses));
    }
    return 0;
}
\end{lstlisting}

\subsection*{A.2 编译与运行命令}
\begin{verbatim}
# 编译
g++ -O2 -std=c++17 -o cache_sim src/cache_sim.cpp

# 运行单个配置（参数：trace 文件、Cache 大小、相联度、块大小）
./cache_sim --trace "trace files/022.li.din" \
            --cache_size 32K --assoc 4 --block_size 64

# 批量运行全部 256 种组合，结果写入 results/results.csv
./run_experiments.sh
\end{verbatim}

\subsection*{A.3 批处理脚本核心：run\_experiments.sh}
\begin{lstlisting}
TRACES=("022.li.din" "047.tomcatv.din" "078.swm256.din" "085.gcc.din")
SIZES=("8K" "16K" "32K" "64K")
ASSOCS=(1 2 4 8)
BLOCKS=(16 32 64 128)
for t in "${TRACES[@]}"; do
  for s in "${SIZES[@]}"; do
    for a in "${ASSOCS[@]}"; do
      for b in "${BLOCKS[@]}"; do
        ./cache_sim --trace "trace files/$t" --cache_size "$s" \
                    --assoc "$a" --block_size "$b" --csv >> results/results.csv
      done
    done
  done
done
\end{lstlisting}

\end{document}
